% =========================================================
% Proof: Mediant Property of Farey Neighbors
% =========================================================

\newpage
\phantomsection
\label{prf:farey-mediant-property}

\begin{remark*}[Return]
\hyperref[thm:farey-mediant-property]{Return to Theorem~\ref*{thm:farey-mediant-property}.}
\end{remark*}

\begin{theorem*}[Mediant property of Farey neighbors]
If $a/b$ and $c/d$ are consecutive fractions in a Farey sequence, then
$(a+c)/(b+d)$ satisfies $a/b < (a+c)/(b+d) < c/d$.
\end{theorem*}

\begin{proof}[Proof (Professional Standard)]
Since $a/b$ and $c/d$ are consecutive Farey neighbors with $b,d > 0$ and
$a/b < c/d$, the \hyperref[thm:mediant-inequality]{Mediant Inequality}
applies directly and gives
\[
\frac{a}{b} < \frac{a+c}{b+d} < \frac{c}{d}.
\qedhere
\]
\end{proof}

\begin{proof}[Proof (Detailed Learning)]
\textbf{Step 1: Verify the hypotheses of the Mediant Inequality.}
We have $a/b < c/d$ and $b, d > 0$ (from the definition of Farey sequence).

\textbf{Step 2: Apply the Mediant Inequality.}
The \hyperref[thm:mediant-inequality]{Mediant Inequality}
(\hyperref[thm:mediant-inequality]{Theorem~\ref*{thm:mediant-inequality}})
states: if $a/b < c/d$ with $b,d > 0$, then
\[
\frac{a}{b} < \frac{a+c}{b+d} < \frac{c}{d}.
\]

\textbf{Step 3: Conclude.}
The hypotheses are satisfied, so the conclusion holds immediately.
\end{proof}

\begin{remark*}[Proof structure]
A direct application of the Mediant Inequality — the Farey context adds
no new content beyond confirming that the hypotheses are met.  The
interest of the result is interpretive: it shows that the mediant is
the natural operation for generating new Farey fractions.
\end{remark*}

\begin{remark*}[Dependencies]
\hyperref[thm:mediant-inequality]{Mediant inequality},
\hyperref[def:farey-sequence]{Farey sequence}.
\end{remark*}
